\documentclass[12pt]{article}
\usepackage[T1]{fontenc}
\usepackage[utf8]{inputenc}
\usepackage{lmodern}
\usepackage{textcomp}
\usepackage{enumitem}
\usepackage{geometry}
\geometry{margin=1in}
\begin{document}
\begin{center}
Answer Key - History - Undergraduate Level Assessment 7 \
\small (Answers are ordered exactly as in the question paper.)
\end{center}

\section*{Multiple Choice}
\begin{enumerate}[label=\arabic*.]
\item \textbf{Answer:} C
\\ \textit{Options:} A) in sub-Saharan Africa, about 10,500 years ago; B) in western Europe, about 3,500 years ago; C) in the Middle East, about 10,500 years ago; D) in North America, about 5,500 years ago; E) in East Asia, about 5,500 years ago
\\ \textit{Note:} The correct answer is supported by archaeological findings that trace the domestication of cattle to the region of the Middle East approximately 10,500 years ago, marking a significant development in early agricultural societies.
\item \textbf{Answer:} A
\\ \textit{Options:} A) mobiliary art.; B) long-distance trade of exotic raw materials.; C) parietal art.; D) a Venus figurine.
\\ \textit{Note:} The "Lion Man" is classified as mobiliary art because it is a portable sculptural piece created during the Upper Paleolithic period, reflecting the artistic expression of that time through small, movable objects rather than fixed site art.
\item \textbf{Answer:} A
\\ \textit{Options:} A) 500,000 years; B) 50,000 years; C) 200,000 years; D) 2,000 years; E) 100,000 years
\\ \textit{Note:} The Pleistocene epoch, characterized by repeated glaciations, is generally understood to have ended around 11,700 years ago, but if considering the broader context of significant climatic shifts leading into the Holocene, the figure of 500,000 years can refer to the timeframe of the Pleistocene's earlier stages, which began approximately 2.6 million years ago.
\item \textbf{Answer:} A
\\ \textit{Options:} A) Mi'kmaq; B) Norse; C) Beothuk; D) Dorset; E) Nootka
\\ \textit{Note:} Genetic evidence indicates that the modern Inuit share ancestry with the Mi'kmaq people, highlighting a historical connection among Indigenous populations in North America, which supports the conclusion that the Inuit are descendants of the Mi'kmaq.
\item \textbf{Answer:} E
\\ \textit{Options:} A) tactical collecting.; B) symbolic foraging.; C) systematic hunting.; D) resource planning.; E) logistical collecting.
\\ \textit{Note:} Logistical collecting refers to a subsistence strategy where groups plan their movements and resource acquisition seasonally to efficiently gather and utilize resources, making it suitable for environments where resource availability fluctuates over time.
\end{enumerate}

\section*{True / False}
\begin{enumerate}[label=\arabic*.]
\setcounter{enumi}{5}
\item \textbf{Answer:} False
\\ \textit{Options:} A) True; B) False
\\ \textit{Note:} The correct answer is: Bolte
\item \textbf{Answer:} False
\\ \textit{Options:} A) True; B) False
\\ \textit{Note:} The correct answer is: Virginia
\item \textbf{Answer:} True
\\ \textit{Options:} A) True; B) False
\\ \textit{Note:} According to the passage: the Valencian Community
\item \textbf{Answer:} True
\\ \textit{Options:} A) True; B) False
\\ \textit{Note:} According to the passage: summaries or numerical lists
\item \textbf{Answer:} False
\\ \textit{Options:} A) True; B) False
\\ \textit{Note:} The correct answer is: 9 th
\end{enumerate}

\section*{Long Form Response}
\begin{enumerate}[label=\arabic*.]
\setcounter{enumi}{10}
\item \textbf{Answer:} eleven
\\ \textit{Note:} Expected answer: eleven
\item \textbf{Answer:} Thomas Cubitt
\\ \textit{Note:} Expected answer: Thomas Cubitt
\end{enumerate}

\vfill
\noindent\textit{End of Answer Key}
\end{document}